\cotefile{136}
\scanpage{267}
Si $\deleted{\unsure{i}}\unsure{\mathcal C|C}\xrightarrow{i}\mathcal M_{N_0}$ est un foncteur exact

\[
\operatorname{Im}\deleted{\unsure{\mathrm{ess}}}\,i\subset\operatorname{Im\,ess}\varphi_{N_0,\infty}
\]

alors on a $i\simeq\varphi_{N_0,\infty}\tau_{\infty,N_0}i$ d'où

\[
(\mathbb R^{\deleted{\unsure{*|n}}}\varphi_{N_0,\infty})(\tau_{\infty,N_0}\circ i)\simeq\mathbb R^*i
\]

en particulier

\[
(R^n\varphi_{N_0,\infty})(\tau_{\infty\,N_0}i(M))=0
\qquad\begin{gathered}\text{si }n>0\\ M\in\mathrm{ob}\,C.\end{gathered}
\]

On est intéressé au = cas particulier qui avant.

\figure{\begin{tikzpicture}\draw (0,0)--(1.45,0);\end{tikzpicture}}

\[
\begin{aligned}
\text{Soit}\qquad &K^\cdot\text{ un complexe de }C\\
&L^\cdot\text{ un complexe de }\mathcal M_{\deleted{\unsure{\infty}}}
\end{aligned}
\]

et supposons donné un hom dans la catégorie dérivée $D(\mathcal M_{N_0})$

\[
i(K^\cdot)\xrightarrow{\alpha}L^\cdot
\]

tel que

\[
\tau_{\infty,N_0}i(K)\xrightarrow[\sim]{\tau_\infty\alpha}\tau_{\infty,N_0}L^\cdot
\]

soit un isom dans $D(\mathcal M_\infty)$, i.e. tel que $\forall n$,

\[
H^n(\alpha):\underset{\substack{\wr\\ i(H^n(K))}}{H^{n}\deleted{\unsure{\varphi}}(i(K))}\longrightarrow H^n(L^\cdot)
\]

soit un $\infty$-isom, i.e. noyau et conoyau \unsure{sont|soit} négligeables. On en conclut

% archive sheet number: 266
\scanpage{268}
% Faint reversed arrow in the upper blank area shows through from the other side.

\mspage{- 83 -}\hfill IX

7.1.3. La projectivité de $\hat B$ signifie, en vertu de EGA III 5.4.5, que l'on peut trouver un Module inversible $\underline{\hat M}$ sur $\hat B$, i.e. une suite de Modules inversibles $\underline M_n$ sur les $B_n$ qui se "recollen\deleted{\unsure{t"}}\unsure{t"} pour $n$ variable, qui soit \uline{ample} relativement à $S$ i.e. qui soit tel que $\underline M_{\unsure{o|0}}$ soit ample. L'idée naturelle serait de partir d'un Module inversible ample $\underline M$ de $G$, et de montrer que pour tout $n$, la restriction $\underline L_n$ de $\underline L$ à $G_n$ provient par image inverse d'un Module inversible $M_n$ sur $B_n$. J'ignore si de tels $M_n$ existent nécessairement (et en doute fort). Nous allons suivre une méthode un peu différente, en construisant une \uline{polarisation} de $B$, i.e. un système projectif de \uline{polarisations} $\rho_n$ des $B_n$ qui se recollent. On sait alors [5] que les doubles de ces polarisations définissent canoniquement des Modules inversibles $\underline M_n$ (symétriques) sur les $B_n$, qui par leur nature canonique se recolleront nécessairement, et fourniront le $\underline M$ cherché.

Nous allons interpréter un élément de $\mathrm{NS}(B_n/S_n)$ comme une correspondance divisorielle symétrique sur $(B_n,B_n)$ [5], i.e. comme une biextension de $(B_n,B_n)$ par $\mathbb G_{m\,S_n}$. En vertu de VIII 3.5, ces biextensions correspondent biunivoquement aux biextensions de $(A_n,A_n)$ par $\mathbb G_{mS}$. Pour obtenir un système cohérent de telles biextensions $W_n$ il suffit de partir d'une

\[
W,\ \text{biextension symétrique de }(A,A)\text{ par }\mathbb G_{mS}\,,\tag*{(7.1.4)}
\]

en prenant les $W_n$ se déduisant de $W$ par les changements de base $S_n\longrightarrow S$. Nous devons de plus supposer que l'élément de $\mathrm{NS}(B_o)$ défini par $W_o$ est une polarisation de $B_o$ pour obtenir ainsi une polarisation de $\hat B$.

% archive sheet number: 267
% stamp: 401 (lower centre, on label); faint partly covered earlier number above.
\scanpage{269}
en appliquant $\mathbb R\varphi_{\unsure{0|N_0},\infty}$

\[
\underbrace{(\mathbb R^*\varphi_{N_0,\infty})(\tau_{\infty,\deleted{\unsure{i}}N_0}i)(K)}
\simeq(\mathbb R^{*}\varphi_{N_0,\infty})(\tau_{\infty N_0}L^\cdot)
\]

\[
\text{i.e.}\qquad i(K)\simeq
\underbrace{(\mathbb R^*\varphi_{N_0,\infty})(\tau_{\infty,N_0}L^\cdot)}_{L'^\cdot}
\]

\deleted{Je dis que l'hom d'augmentation}

Mais c'est là un isom dans $D(\mathcal M_{N_0})$ \unsure{seult|seulement}, -- rien de plus précis que la donnée de départ -- on aimerait un isom dans $D(C)$. On va supposer $i$ pl.\ fidèle et avec un adjoint à droite $k$,

\[
\text{donc}\qquad k\circ i\simeq\mathrm{id}_C
\]

et on va supposer $k$ \uline{exact}. \unsure{Appliquant|Appliquons} $k$ à l'isom obtenu, on trouve

\[
K^\cdot\simeq k\bigl(\mathbb R^*\varphi_{N_0,\infty}(\tau_{\infty,N_0}L)\bigr)
\]

\inserted{$\simeq(\mathbb R^*(k\varphi_{N_0\infty}))(\tau_{\infty N_0}L^\cdot)$}

i.e. $K^\cdot$ se \unsure{reconstitue} dans $D(\unsure{\mathcal C|C})$ à l'aide de $\tau_{\infty,N_0}(i(K^\cdot))$ \deleted{\unsure{dans}} $\tau_{\infty,N_0}(L^\cdot)$ dans $D(\mathcal M_\infty)$.

\deleted{On va simplifier l'expression obtenue, en considérant}

\[
\deleted{\varphi_{N_0,\infty}(\tau_{\infty,N_0}(L^\cdot))\longrightarrow\mathbb R^*\varphi_{N_0,\infty}(\tau_{\infty N_0}L))}
\]

\deleted{\uline{Lemme} Cet hom est un quasi-isomorphisme}

\deleted{\uline{\unsure{Corollaire}} $\begin{aligned} K^\cdot&\simeq k\varphi_{N_0,\infty}(\tau_{\infty,N_0}L^\cdot)\\ &\simeq k\varphi_{N_0,\infty}(\tau_{\infty N_0}iK).\end{aligned}$}

% archive sheet number: 268
\scanpage{270}
\mspage{- 82 -}\hfill IX

par un tore. Nous avons posé comme d'habitude $A_n=A\times_S S_n$, où $S_n=\operatorname{Spec}(V/\underline m^{n+1})$. Nous nous intéresserons surtout au cas où $S$ est un trait, et où la fibre générique $A_K$ est un schéma abélien à réduction semi-stable sur $S$, $A$ étant le modèle de Néron de $A_K$.

Soit $T_{\unsure{o|0}}$ le tore maximal de $A_o$, de sorte que $B_o=A_o^o/T_o$ est un schéma abélien. La construction de 5.1 s'applique alors pour fournir un sous-tore formel $\hat T$ du complété formel $\hat A$. On sait que pour tout $n$, le quotient

\[
B_n=A_n^o/T_n\tag*{(7.1.1)}
\]

est représentable (SGA 3 $\mathrm{VI}_A$ 3.2), plus précisément c'est un schéma plat et de type fini sur $S_n$, dont l'image inverse sur $S_o$ est le schéma abélien $B_o$. Par suite $B_n$ est un \uline{schéma abélien} sur $S_n$, et pour $n$ variable, les $B_n$ se déduisent les uns des autres par changement de base, donc ils définissent un \uline{schéma abélien formel}

\[
\hat B=(B_n)_{n\geq0}\,.\tag*{(7.1.2)}
\]

Nous allons donner des conditions impliquant que ce schéma abélien est \uline{algébrisable} (EGA III 5.4.1), plus précisément que c'est le complété formel d'un schéma abélien projectif sur $S$, que nous noterons alors $B$. En vertu de loc.cit. $B$ est déterminé à isomorphisme unique près (comme $S$-schéma ; la structure de groupe de $\hat B$ provient alors de façon unique d'une structure de groupe de $B$, en vertu de loc. cit. également). Ce schéma abélien $B$ sera alors appelé \uline{la partie abélienne} du schéma $A$, et parfois noté $A^{\mathrm{ab}}$.

% archive sheet number: 269
% stamp: 400 (lower centre label; the pencilled 400 above it is also visible).
\scanpage{271}
\diagram{\begin{tikzcd}[ampersand replacement=\&,column sep=huge]
 C \arrow[r,"{\deleted{\unsure{\varphi}}\tau_{\infty N_0}i=\sigma}"'] \&
 \mathcal M_\infty \arrow[l,bend left=65,"{k\varphi_{N_0\infty}=\psi}"']
\end{tikzcd}}

$\sigma$ pl.\ fidèle, adj.\ à droite $\psi$

\[
\psi\sigma\xleftarrow{\sim}\mathrm{id}_C\qquad
\left|\begin{array}{l}\deleted{\psi\text{ exact}}\\ \sigma\text{ exact}\\ k\text{ exact}\end{array}\right.
\]

mais $\psi$ peut-être pas exact

\[
\deleted{\mathbb R}\quad\sigma(K^\cdot)\longrightarrow L_\infty^\cdot\qquad\text{quasi-isom dans }\mathcal M_\infty
\]

\diagram{\begin{tikzcd}[ampersand replacement=\&,column sep=small,row sep=small]
 \deleted{(R^{\unsure{*|\alpha}}\psi)\underbrace{\sigma(K^\cdot)}_{\substack{\Vert\\R^*\sigma K}}}
 \arrow[r,"\sim"] \arrow[d,equals] \&
 \deleted{R^*\psi(L_\infty^\cdot)} \&
 \deleted{\psi(L_\infty^\cdot)}\arrow[l,"{\overset{?}{\sim}}"']\\
 \deleted{R^*(\underbrace{\psi\sigma}_{\unsure{\text{est id}|\text{exact}|\simeq\mathrm{id}}})(K)\simeq K}\&
 \deleted{\deleted{\unsure{\mathbb R^*\psi|R^*\psi}}\,H^*(\mathbb R\psi(L_\infty^\cdot))}\&\\
 \&\deleted{E_2^{pq}=R^q\psi(H^p(L_\infty^\cdot))}\arrow[u,Rightarrow]\&
\end{tikzcd}}

\deleted{si $R^q\psi(H^{\unsure{p}}(L_\infty^\cdot))=0$ pr $q>0$}

\[
\begin{aligned}
&\deleted{E_1^{p,q}\mathbin{\unsure{-|\to}} E_1^{p+1,q}}\\
&\deleted{E_2^{p,q}\mathbin{-} E^{p+2,q-1}}
\end{aligned}
\]

\deleted{On trouve}

\[
\deleted{H^n(\mathbb R^{\unsure{*}}\psi(L_\infty^\cdot))\simeq\psi(H^n(L_\infty^\cdot))}
\]

\[
\begin{aligned}
&\deleted{i(K^\cdot)\longleftarrow L^\cdot}\\
&\deleted{\tau_{\infty N_0}i(K^\cdot)\xleftarrow{\sim}\tau_{\infty N_0}(L^\cdot)}\\
&\deleted{K^\cdot\simeq\mathbb R^*(k_0\varphi_\infty)(\tau_{\infty N_0}(L\unsure{.|{}^{\cdot}}))}\\
&\deleted{D(\unsure{\mathrm{Mod}(W)|\mathrm{Mod}(k)})\xrightarrow[\text{pl. fidèle}]{}D(\mathcal M_\infty)}
\end{aligned}
\]

% archive sheet number: 270
\scanpage{272}
\mspage{- 81 -}\hfill IX

opère (de façon non nécessairement fidèle) sur $k'$ (le corps des invariants étant $k$), et sur $Y_{k'}$ da façon compatible avec son action sur $k'$. Or, d'après la théorie de DIEUDONNÉ [16 bis], au $p$-pro-groupe de Barsotti-Tate $Y_{k'}$ est canoniquement associé un module $M$ libre de type fini (de rang égal au rang de $Y_{k'}$) sur l'anneau $W(k')$ des vecteurs de Witt sur $k'$, et $G$ opère évidemment sur $M$ de façon compatible avec ses opérations sur $W(k')$ (via ses opérations sur $k'$). On en déduit une opération $W(k')$-linéaire du sous-groupe d'inertie $\Gamma$ de $G$ sur $M$ (qui est un quotient du groupe d'inertie $I$), d'où une fonction trace sur $\Gamma$ donc sur $I$, qui est la fonction cherchée.Elle est a priori à valeurs dans $W(k')$ (*), mais coïncide avec la fonction analogue définie en termes des $T_\ell$ pour $\ell\ne p$, lorsque l'opération de $G$ sur $Y_{S'}$ s'obtient à partir d'une opération de $G$ sur un schéma en groupes lisse et $p$-divisible $C'$ sur $S'$, en prenant $Y=T_p(C')$ (cf. [25 bis, th.7 p.737]) ; elle est donc dans ce cas à valeurs entières. On procède de façon analogue dans le cas de 4.3 b) pour définir une fonction trace sur $\pi^!$, en utilisant le procédé de la démonstration de 4.3.2 pour "linéariser" l'opération dans $M$ définie par un élément $g\in\pi$ qui se trouve au-dessus d'un $f_q^n$ avec $n\geq0$.

7. \uline{L'extension de Raynaud} $G^{\handwritten{\unsure{\natural|\sharp}}}$ \uline{sur} $S$ \uline{attachée au schéma abélien} $A_K$ \uline{à réduction semi-stable}. \uline{Dualité des schémas abéliens} $B,B'$ \uline{sur} $S$.

7.1. Soit $S$ un schéma noethérien spectre d'un anneau $V$, séparé et complet pour la topologie définie par un idéal\unsure{.} $\underline m$ de $V$. Soit $A$ un schéma en groupes lisse sur $S$, $A^{\unsure{o|0}}$ sa composante neutre, et supposons que $A_o/A_o^o$ soit fini (donc fini étale) sur $S_o$, et que $A_o^o$ soit une extension d'un schéma abélien

\noindent (*) En fait, il est immédiat que cette fonction est toujours à valeurs dans $W(k)$, étant invariante par $\operatorname{Gal}(k'/k)$.

% Heavy reversed arrows at « neutre » and « et que » match the two left arrows on scan 271.
% archive sheet number: 271
% stamp: 399 (lower centre label; earlier 399 partly covered).
\scanpage{273}
\deleted{$\psi\sigma=\mathrm{id}$}

\diagram{\begin{tikzcd}[ampersand replacement=\&,column sep=huge]
 D^+(C)\arrow[r,"{\unsure{\mathbb R|R}\sigma}"]\&
 D^+(\mathcal M_\infty\arrow[l,bend left=40,"{R\psi}"']
\end{tikzcd}}

\margin{$\begin{aligned}\psi\sigma&\xleftarrow{\sim}\mathrm{id}_C\\
\sigma\psi&\mathbin{\unsure{\to|-}}\mathrm{id}_{\mathcal M_\infty}\end{aligned}$}

\diagram{\begin{tikzcd}[ampersand replacement=\&,column sep=small,row sep=large]
 \&\operatorname{Hom}(R\sigma K,\underbrace{R\sigma R}_{R(\sigma\psi)(L)}\psi(L)\arrow[dl]\&\\
 \operatorname{Hom}(R\sigma(K),L)\arrow[rr,"{\overset{?}{\sim}}"]\arrow[dr]\&
 \&\operatorname{Hom}(K,R\psi(L)\arrow[ul]\\
 \&\operatorname{Hom}(\underbrace{R\psi R\sigma K}_{R(\psi\sigma)(K)\simeq K},R\psi L)
 \arrow[ur,phantom,"{\simeq}"{sloped}]\&
\end{tikzcd}}

\diagram{\begin{tikzcd}[ampersand replacement=\&]
 R\sigma R\psi R\sigma K,\arrow[d]\& R\sigma R\psi L\arrow[d]\\
 R\sigma K\&L
\end{tikzcd}}

\[
D^+C\xrightarrow[\substack{\text{pl fid.}\\\text{compatible avec}\\\text{multiplication}}]{}D^+\mathcal M_\infty
\]

% archive sheet number: 272
\scanpage{274}
\mspage{- 80 -}\hfill IX

6.4. Notons que lorsque $\ell\ne p$, nous avons mis en évidence dans 4.3 a) une fonction à valeurs entières remarquable sur le groupe d'inertie $I$ de $S$, d'ailleurs indépendante du choix de $\ell$, et dont la connaissance, dans le cas d'un corps résiduel parfait, suffit à reconstruire l'exposant conducteur de $A$ (4.5), en prenant le produit scalaire, sur un quotient assez grand $\Gamma=I/I'$ de $I$, de la fonction précédente et du conducteur d'Artin sur $\Gamma$ (cf. (4.6.2)). Il peut être intéressant d'observer que cette fonction à valeurs entières sur $I$ peut se réconstituer également par la connaissance de $T_\ell(A_K)$ pour $\ell=p$ - du moins (pour le moment) en inégales caractéristiques. Lorsque le corps résiduel $k$ est fini, une observation analogue s'applique à la fonction à valeurs dans $\mathbb Q$, définie sur le sous-groupe $\pi^!$ de $\pi$ formé des éléments dont l'image dans $\pi_{\unsure{o|0}}$ est de la forme $f_q^n$ ($n\in\mathbb Z$), considérée dans 4.3 b). Indiquons le principe de la vérification de ces assertions. On utilise la filtration (6.3.1), et on est réduit à attacher une fonction convenable (sur $I$ resp. sur $\pi^!$) à un $p$-pro-groupe de Barsoti-Tate $Y$ sur $K$ qui a "essentiellement bonne réduction sur $S$ ", au sens précisé dans 6.5. ($Y$ sera un des trois quotients provenant de la filtration (6.3.1).) Dans le cas de 4.3 a), on se \unsure{ramène|remène} au cas d'un corps résiduel parfait (ou même algébriquement clos, si on veut) par extension résiduelle convenable dans $S$ (EGA $0_{\mathrm{III}}$ 10.3.1). Prenons alors une extension galoisienne finie $K'$ de $K$ telle que $Y_{K'}$ ait bonne réduction sur $S'$, d'où un groupe de Barsotti-Tate $Y_{S'}$ sur $S'$ et une opération de $G$ sur $Y_{S'}$ compatible avec son action sur $S'$. Si $k'$ est le corps résiduel de $S'$, $G$

% Reverse impression from scan 273: the backward R in the upper margin and the marks across « on se ramène au cas ».
% archive sheet number: 273
% stamp: 398 (lower centre label; the pencilled 398 above is also visible).
\scanpage{275}
\[
\begin{aligned} C&\xrightarrow{i}\mathcal M\\ M&\longmapsto M\otimes_k S\end{aligned}
\]

\[
i\xrightarrow{\sim}\varphi_{0,N}\tau_{N,0}\deleted{\unsure{i}}i
\]

\diagram{\begin{tikzcd}[ampersand replacement=\&,column sep=large]
 (x_h)_{h\geq N}\arrow[r,shift left=.7ex,"{0\leq i\leq N}"]\arrow[r,shift right=.7ex]\&
 c_{\ell,\unsure{h|k}}x_h-c_{\ell,h-i}x_{h+\ell}
\end{tikzcd}}

\diagram{\begin{tikzcd}[ampersand replacement=\&,column sep=large]
 M\arrow[r,hook]\& M^{[N,+\infty[}\arrow[r,shift left=.6ex]\arrow[r,shift right=.6ex,squiggly]\&
 M^{[N,+\infty]\times\mathbb N}\\
 x\arrow[r,mapsto]\&(c_{i,h-i}\,x)_{h\geq N}\&
\end{tikzcd}}

\[
M\longmapsto\underline{\operatorname{Hom}}_k^*(S,M)
\]

\[
\begin{aligned}
\varphi^{0,N}\tau^{N,0}\circ i&\longrightarrow i\qquad
\text{syst. des }\underline{\operatorname{Hom}}_k(S_i,M)\\
&\mathrm{Ab}\qquad\mathrm{Ind}\,(\mathrm{Ab})\qquad\simeq S\otimes_k M
\end{aligned}
\]

\[
M\mathbin{\unsure{\hat\otimes|\otimes}}_k S_i\mathbin{\to}\underline{\operatorname{Hom}}^i(\tau_{N-i}(S),M\otimes_k S)
\]

$C$ Cat. ab. $k$-linéaire (i.e. $k\longrightarrow\operatorname{End}(\mathrm{id}_C)$)

$X$ objet de $C$

$M$ un $k$-module à g (resp à \unsure{droit|droite}

\[
\begin{aligned}
&\underline{\operatorname{Hom}}(M,X)\in C\qquad\text{défini par}\\
&\operatorname{Hom}(Y,\underline{\operatorname{Hom}}(M,X))=\operatorname{Hom}_k(M,\operatorname{Hom}_C(Y,X))
\end{aligned}
\]

\[
\begin{aligned}
&M\otimes_k X\qquad\text{défini par}\\
&\operatorname{Hom}(M\otimes_k X,Y)\simeq\operatorname{Hom}_k(M,\operatorname{Hom}_C(X,Y))
\end{aligned}
\]

Dans la mesure où ça existe $M\longmapsto\underline{\operatorname{Hom}}(M,\unsure{N|X})$ et $M\deleted{\unsure{\mapsto}}\mapsto M\otimes_k X$ est exact à droite --

% archive sheet number: 275
\scanpage{276}
\mspage{- 127 -}\hfill IX

10.5. \uline{Réduction de} 10.4 \uline{à un cas particulier}.

10.5.1. L'unicité d'un $u$ satisfaisant la condition énoncée (même pour un seul $\ell$) est claire. D'autre part, l'existence signifie simplement que chaque $u_\ell$ applique en fait $\underline M\otimes\underline M'$ dans $\unsure{\mathbb Z_S}$, et que l'homomorphisme ainsi obtenu $\underline M\otimes\underline M'\longrightarrow\mathbb Z_S$ est indépendant de $\ell$. En fait, il suffit même de prouver l'assertion analogue pour les $u_\ell\otimes_{\mathbb Z_\ell}\mathbb Q_\ell$ et un $u_{\mathbb Q}:\underline M_{\mathbb Q}\otimes\underline M'_{\mathbb Q}\longrightarrow\mathbb Q_S$.

Supposons alors qu'on ait établi l'existence après changement de base $S'\longrightarrow S$ comme dans 9.3. Alors il résulte immédiatement de (10.3.5) qu'il existe un homomorphisme

\[
u:\underline M\otimes\underline M'\longrightarrow\frac{1}{e(S'/S)}\,\mathbb Z_S\subset Q_S
\]

compatible avec les $u_\ell\otimes_{\mathbb Z_\ell}\mathbb Q_\ell$, d'où a). Quant à b), il est trivial alors que s'il est vrai pour $A_{K'}$, il l'est aussi pour $A_K$.

Par suite, nous pouvons, pour prouver 10.4, nous permettre n'importe quel changement de traits $S'\longrightarrow S$ dominant. Ceci nous permet en particulier de supposer $S$ complet à corps résiduel algébriquement clos. Dans ce cas, $\underline M$ et $\underline M'$ sont des schémas en groupes constants, de valeur $M,M'$ des modules libres de type fini sur $\mathbb Z$, et a) et b) se reformulent avantageusement en termes de ces derniers.

10.5.2. Supposons maintenant que $A_K$ soit isomorphe à un facteur direct d'un schéma abélien $\bar A_K$ sur $k$ à réduction semi-stable, et que 10.4 soit prouvé pour ce dernier. Alors il en résulte aussitôt que 10.4 est valable également pour $A_K$, compte tenu du comportement fonctoriel évident des réseaux $\underline M,\underline M'$, et des accouplements de monodromie, vis-à-vis des sommes

% archive sheet number: 276
% stamp: 445 (lower centre label; 445 in pencil above it).
% Reversed strokes from scan 275 show through, including the horizontal after the tensor product in line 11.
\scanpage{277}
\[
P\mathbin{\unsure{\to|-}} Q-R\qquad P\to Q-Z'-0\qquad Z'-R
\]

\diagram{\begin{tikzcd}[ampersand replacement=\&,column sep=small,row sep=small]
 \deleted{P\to Z\to0}\arrow[d]\&\deleted{0\to Z\to Q\to Z'\to0}\\
 \text{conservé par }T\&0-Z'-R
\end{tikzcd}}

\[
T(P)\mathbin{-}T(Q)\mathbin{-}T(Z')\to0
\]

\diagram{\begin{tikzcd}[ampersand replacement=\&,column sep=small,row sep=small]
 T(Z)\arrow[r]\&T(Q)\&
 T(Q)/\operatorname{Im}T(P)\simeq T(Z')\arrow[d,hook,"?"]\\
 T(P)\arrow[u]\&\&T(R)
\end{tikzcd}}

\[
0\to Z\mathbin{-}Q\to Z'\subset R
\]

\[
T(Z)\to T(Q)\mathbin{-}T(Z')\to0
\]

$I\subset R$

\margin{\unsure{reste une injection} en tensorisant\\
\unsure{transformé en surjection|transforme en surjection|transformer en surjection} par Hom}

\[
Z'\hookrightarrow R
\]

\diagram{\begin{tikzcd}[ampersand replacement=\&,column sep=small,row sep=small]
 P\arrow[r,dash]\&Q\arrow[r,two heads]\& I\arrow[r,hook]\arrow[d]\&R\\
 \&\&0\&
\end{tikzcd}}

\[
\begin{aligned}
&T(P)\leftarrow T(Q)\longleftarrow T^{\unsure{0|\circ|\bullet}}(I)\leftarrow0\\
&T(I)=\operatorname{Ker}(T(Q)\mathbin{\unsure{-|\to}}T(\unsure{P}))
\end{aligned}
\]

\[
\begin{aligned}
&\operatorname{Hom}(P,M)\\
&\operatorname{Hom}_{C,\unsure{S|s}}(N,\operatorname{Hom}(P,M))\\
&\deleted{\unsure{24\,24}}\\
&\deleted{\operatorname{Hom}(N\otimes P,M)}\\
&\operatorname{Hom}_{\deleted{\unsure{\otimes}}}
\end{aligned}
\]

\[
T(\unsure{R|Q})
\]

\diagram{\begin{tikzcd}[ampersand replacement=\&,row sep=small,column sep=small]
 \mathbb Z^{[N,+\infty[\times\mathbb N}
 \&\mathbb Z^{[N,+\infty]}\&\mathbb Z\&\\
 P\arrow[u,equals]\arrow[r,dash]\& Q\arrow[u,equals]\arrow[r]\arrow[dr]\&
 R\arrow[u,equals]\arrow[r]\&0\\
 \&\&R'\arrow[u,phantom,"{\bigcup}" description]\&
\end{tikzcd}}

\[
(\lambda_{h,\ell})\underset{\substack{h\geq N\\\ell\geq0}}{\longrightarrow}
\sum_{\substack{\deleted{\unsure{h,\ell}}\\h,\ell}}\lambda_{h,\ell}
\bigl(\deleted{\unsure{c_{h}}}c_{\unsure{\ell h|ih|\ell,k}}e_h-c_{\ell,h-i}\,e_{h+\ell}\bigr)
\]

\[
\sum\ (\lambda_h)_{h\geq0}\quad\searrow\quad
\deleted{\sum c_{i,h-i}\,\unsure{\lambda_h|x_h|\lambda_k}}
\]

% archive sheet number: 277
\scanpage{278}
\mspage{- 108 -}\hfill IX

où les flèches verticales désignent les foncteurs évidents induits par le foncteur exact $r$, et où $\omega'$ est la flèche analogue à $\omega$ (9.3.5), avec $C$ remplacé par $C'$. On en déduit des compatibilités analogues des structures envisagées dans 9.3.8 avec le foncteur $r$, que nous laissons au soin du lecteur à expliciter.

Nous utiliserons la conséquence suivante de cette commutativité. Supposons que $\operatorname{EXTPAN}(E,F)$ ne soit pas vide, et choisissons alors un objet $X$ dans $\operatorname{EXTPAN}(E,F)$. Alors grâce à $G'\longmapsto G'\wedge X'$ (où $X'=r(X)$), la catégorie $\operatorname{EXTPAN}(E',F')$ est équivalente à la catégorie $\operatorname{EXT}(Q',P')$, donc un élément $Y'$ de $\operatorname{EXTPAN}(E',F')$ est connu à isomorphisme près quand on connaît la classe d'isomorphie de l'élément correspondant de $\operatorname{EXT}(Q',P')$, i.e. quand on connaît l'élément correspondant de $\operatorname{Ext}^1(Q',P')$, soit

\[
c_X(Y')\in\operatorname{Ext}^1(Q',P')\,.
\]

Lorsqu'on remplace $X$ par un autre objet $X_1=H\wedge X$, où $H$ est une extension de $Q$ par $P$, on voit aussitôt que l'on obtient

\[
c_{X_1}(Y')=c_X(Y)-\mathrm{cl}(H')\qquad,
\]

grâce à la formule d'associativité (9.3.7). Par suite, pour $Y'$ fixé, et $X$ variable, les $c_X(Y')$ parcourent exactement une classe

\[
c(Y')\in\operatorname{Ext}^1(Q',P')/\operatorname{Im}\operatorname{Ext}^1(Q,P)\qquad,\tag*{(9.3.13)}
\]

qui est donc canoniquement associée à l'extension panachée $Y'$ de $E'$ par $F'$. Par construction, sa nullité est nécessaire et suffisante pour que l'extension panachée $Y'$ soit isomorphe à la transformée par $r$ d'une extension panachée $X$ de $E$ par $F$.

% Typescript upside down in the archive image.
% archive sheet number: 278
% stamp: 426 (number in pencil and on label).
\scanpage{279}
\diagram{\begin{tikzcd}[ampersand replacement=\&,row sep=small]
 s=k[u]\arrow[d,bend right=20]\\
 u^n=u_n\text{ de degré }-n
\end{tikzcd}}

un $\unsure{s^k|S^k}=k\{\{T\}\}$-module à gauche un $S$-module (\unsure{gradué|gradués}) avec $S=k\{T\}$

\[
M\in\operatorname{Mod}(k)
\]

\[
\begin{aligned}
&M\longrightarrow M\otimes_k s\qquad\text{un foncteur}\\
&\underset{\substack{\Vert\\ C}}{\operatorname{Mod}(k)}
 \xrightarrow{i_*}\operatorname{Mod}\,\mathrm{grad}\ k\{T\}=\mathcal M
\end{aligned}
\]

si $N\in\mathbb Z$ / $\sigma_N(P)=P/\tau_{N+1}(P)$ "cotronqué" \inserted{ou "tronqué à gauche"} en dim $N$

\[
\sigma_N(P)_i=\begin{cases}P_i\text{ si }i\leq N\\0\text{ si }i>N\end{cases}
\]

Considérons le composé $\sigma_N\circ i_*$ ($N\leq0$, on fer. $N\mathbin{\unsure{\mapsto|\to}}-\infty$)

\uline{Th}. Le foncteur $\tau_N i_*$ est pl. fid.

\[
\operatorname{Hom}_{\mathcal M}\deleted{\unsure{\mathrm{grad}|\mathrm{gr}}}(M\otimes_k s,\tau^{\deleted{\unsure{N}}}\unsure{\mathcal N|N|P})
=\operatorname{Hom}_k(M,\operatorname{Hom}_{\mathcal M}(s,\mathcal N))
\]

\[
\begin{aligned}
&\operatorname{Hom}_{\mathcal M_N}(\tau^{N}\deleted{\unsure{N}}M\otimes_k s,\tau^N N\otimes_k s)
 =\operatorname{Hom}_{\mathcal M}(M\otimes_k s,i^{N,\infty}(N\otimes_k s))\\
&\hspace{3em}=\operatorname{Hom}_k(M,\underbrace{\operatorname{Hom}_{\mathcal M}(s,i^{N,\unsure{*}}\tau^N\mathord{-})}_{\unsure{N_0}})
\end{aligned}
\]

\[
\begin{aligned}
&\operatorname{Hom}_{\mathcal M}(s,\tau^N N\otimes_k s)\simeq N_0\quad??\\
&\qquad\Vert\ \operatorname{Hom}_{V(s),\mathrm{grad}}(V(N\otimes s),V(s))
\end{aligned}
\]

\[
u_{n+h}\xrightarrow[T^{(h)}]{c_{n,h}}u_n\quad/\quad{P_\infty}
\qquad 1
\]

\[
\begin{aligned}
&\operatorname{Hom}_S(s,M\otimes_k S)\mathord{-}\\
&\operatorname{Hom}_{\mathrm{coalg}}(s,M\otimes_k S)\\
&=\operatorname{Hom}_{\unsure{V(s)\text{-mod}}}(\operatorname{Hom}_k(M,V(s)),V(s))\\
&\deleted{\unsure{D}}\,V_S(M)=V(M\otimes_k s)\\
&\qquad=(k'\mathbin{\unsure{-|\to|\mapsto}}\operatorname{Hom}_k(M,s^{k'}))\\
& V(s)\text{ un dual}
\end{aligned}
\]

\[
\begin{aligned}
&\unsure{y_n|\eta_n}\in P_{-n}\qquad \deleted{\unsure{\xi_n}}\,T^{(h)}y_{n+h}=c_{n,h}y_n\\
&\deleted{\unsure{y_{n+h}\,u_{n+h}=c}}\\
&y_{n\deleted{\unsure{+h|+k}}}=\xi_n u_n.\\
&\xi_{n+h}\underbrace{T^{(h)}u_{n+h}}_{c_{n,h}u_n}=\xi_n c_{n,h}u_n\\
&\hspace{5em}\operatorname{Hom}\ V_S(M)\\
&c_{n,h}(\xi_{n+h}-\xi_n)=0
\end{aligned}
\]

% archive sheet number: 279
\scanpage{280}
\mspage{- 33 -}\hfill IX

de traits $S'\longrightarrow S$. Enfin, si (iv) est satisfaite pour $A_K$, elle l'est aussi manifestement pour tout sous-schéma abélien et tout schéma abélien quotient de $A_K$, et en particulier pour tout schéma abélien \uline{isogène} à $A$ (par exemple pour $A'_K$).

\uline{Corollaire} 3.5.2. \uline{Soit} $g$ \uline{un générateur topologique de} $\mathbb Z_\ell(1)=T_\ell(\bar k\unsure{\mathord{\ast}|{}^{*}})$. \uline{Les conditions équivalentes de} 3.5 \uline{équivalent aussi à la condition suivante} :

(v) \uline{L'action du groupe d'inertie} $I$ \uline{sur} $U=T_\ell(A(\bar K))$ \uline{se fait à travers le plus grand quotient pro}-$\ell$-\uline{fini} $I(\ell)$ \uline{de} $I$, \uline{qui est canoniquement isomorphe à} $\mathbb Z_\ell(1)$ (I\hspace{2.5em}), \uline{et l'opération} $g_U$ \uline{du générateur} $g$ \uline{sur} $U$ \uline{satisfait la condition}

\[
(1-g_U)^2=0\quad.\tag*{(3.5.3)}
\]

En effet, pour toute action unipotente continue d'un groupe profini $I$ sur un espace vectoriel $\ell$-adique, on voit aussitôt que l'image de $I$ est un pro-$\ell$-groupe. Il s'ensuit que (v) équivaut à (iv).

Nous pouvons maintenant prouver le résultat central du présent exposé :

\uline{Théorème} 3.6. (\uline{théorème de réduction semi-stable}). \uline{Soient} $S$ \uline{un schéma noethérien régulier connexe de dimension} 1, $K$ \uline{le corps des fonctions rationnelles de} $S$, $A_K$ \uline{un schéma abélien sur} $K$. \uline{Alors il existe une extension finie galoisienne} $K'$ \uline{de} $K$, \uline{telle que le schéma abélien} $A_{K'}=A_K\otimes_K K'$ \uline{ait réduction semi-stable} (3.4) \uline{sur le normalisé} $S'$ \uline{de} $S$ \uline{dans} $K'$.

% archive sheet number: 280
% stamp: 351 (number in pencil and on label).
\scanpage{281}
\[
\begin{aligned}
&M\otimes s\otimes M'\otimes s&&M\otimes M'\otimes s\otimes s\\
&\mathrm{id}_M\otimes u_s\otimes\mathrm{id}_{M'}\otimes\mathrm{id}_s
 &&\mathrm{id}_{M'}\otimes u_s\otimes\mathrm{id}_s\\
&\mathrm{id}_M\otimes\mathrm{id}_s\otimes\mathrm{id}_{M'}\otimes u_s
 &&\mathrm{id}_{M'}\otimes\mathrm{id}_s\otimes u_s
\end{aligned}
\]

\[
\begin{aligned}
&s\boxtimes s=s\\
&s\boxtimes C\simeq C\qquad s\mathbin{\unsure{\perp|{\llcorner}|L}}\\
&\boxed{(M\otimes_k s)\boxtimes_s C\simeq M\otimes_k C}\quad?
\end{aligned}
\]

\[
\begin{aligned}
&\bigl(M\otimes_k\overline{s)\otimes_k C}\ \longleftarrow\ M\otimes_k C\\
&\mathrm{id}_M\otimes u_s\otimes\mathrm{id}_C\\
&\mathrm{id}_M\otimes\mathrm{id}_s\otimes u_C
\end{aligned}
\]

\[
\begin{aligned}
&M\ \unsure{\circ|O}\ N\\
&\qquad\odot
\end{aligned}
\]

\diagram{\begin{tikzcd}[ampersand replacement=\&]
 C\arrow[r]\&C\otimes_k s\arrow[r,shift left]\arrow[r,shift right]\&C\otimes_k s\otimes_k s
\end{tikzcd}}

\deleted{\[
V(s)\otimes V(C)\mathbin{\unsure{-|\to}}V(C')\cdot\unsure{V(C'}
\]}

\[
V(M\otimes_k s)(\unsure{k'|k})=\operatorname{Hom}_k(M.\qquad
\operatorname{Hom}_k(M,s^{k'})\times C^{k'}\longrightarrow C'^{k'}
\]

\deleted{$\unsure{M_k}$}

$\operatorname{Hom}_{k'}(M\otimes_k s,k')=\operatorname{Hom}_k(M,s^{k'})$ est un $\unsure{e^k|s^k|S^{k'}}$-module !

% The following block is upside down relative to the preceding lines.

\deleted{\[
\begin{aligned}
&\alpha_*(\unsure{1|A|\mathbb{1}}),\,\gamma_0^{(i)}
 &&\bigl(A^*(\Sigma,\Sigma_+,t)\bigr)^{k'}\\
&\Delta\gamma_0^i=\sum_{j+h=i}\gamma_0^j\otimes\gamma_0^h
 &&\quad=A_*\!^{\unsure{*|n}}(\Sigma^{k'},\Sigma_+^{k'},t^{k'})\\
&\beta_{\unsure{*|\lambda}}(t)=\sum t^{(i)}\gamma_0^i=\exp(t\gamma_0)
 && A_*^{\unsure{n|\times}}(S,S^+,t)\\
&\unsure{\mathfrak s|s}\quad\deleted{\unsure{\mathscr A}}\quad A^{n}\hat\otimes A^n\longleftarrow A^n
 && A^*(\Sigma,\Sigma_+,t)^\vee\\
&\qquad\Sigma/k &&\qquad\Vert\\
&&&A_*^{\wedge}(\Sigma^\vee,\Sigma_+^\vee,t^\vee)
\end{aligned}
\]}

\deleted{\figure{\begin{tikzpicture}
\node[draw,rounded corners=19pt,inner sep=5pt] (draft) {$
\deleted{\unsure{P}}\quad
\begin{aligned}
&\Sigma\\
&\mathbin{\unsure{\cup|\supset}}\\
&\deleted{\unsure{K}}\\
&\Sigma'
\end{aligned}
\quad\bigg|\quad
\begin{aligned}
&k\text{-coalgèbre }\deleted{\unsure{\text{graduée}}}\ \inserted{\unsure{\text{divisée}}}\ \text{ avec sous-coalgèbre}\\
&\text{et }\unsure{\text{copuiss.}|\text{à puiss.}|\text{augment.}|\text{puiss.}}\ \unsure{\text{divisées}|\text{divisée}}\ \deleted{\unsure{\text{sur}}}\ \unsure{\text{unit}|\text{unit.}}\ \unsure{\text{celle-ci.}|\text{celle-c.}}\\
&\text{enfin on demande }\Sigma\xrightarrow{t}k\\
&\text{nulle sur }\Sigma'
\end{aligned}
$};
\draw (draft.north west)--(draft.south east);
\end{tikzpicture}}}

\[
\begin{aligned}
&\quad\ \ t_k^{\unsure{n}}\\
&0\longrightarrow\boxed{\unsure{\mathbb Z|Z|\Sigma'}}^{k+}\longrightarrow\underbrace{\Sigma^k}_{\text{anneau}}\mathbin{\unsure{-|\to}}\underset{\text{anneau}}{\Sigma'{}^k}
\qquad /\quad k\{\{T\}\}\{\{X_0,-X_n\}\}\\
&\qquad\text{idéal}
\end{aligned}
\]

% archive sheet number: 281
